y view used in
Section~\ref{subsec:distance_dependent_utility}, let
\(\widehat J_{\mathcal B}(\theta)\) denote an empirical estimate of task
performance on an evaluation batch \(\mathcal B\). For a negative sample
\(z=(s,a)\), define its raw repulsive update as
\begin{equation}
    g_\theta^-(z)
    =
    -\widehat A^-(z)
    \nabla_\theta \log \pi_\theta(a\mid s).
    \label{eq:raw_negative_update}
\end{equation}
Assigning it a weight \(w\in[0,1]\) produces the parameter step
\begin{equation}
    \Delta\theta(z;w)
    =
    \eta w g_\theta^-(z).
    \label{eq:weighted_negative_step}
\end{equation}
We define the empirical utility of this weighted negative update as
\begin{equation}
    \widehat{\mathcal U}_{\mathcal B}
    (z;w\mid\theta)
    =
    \widehat J_{\mathcal B}
    \!\left(
        \theta+\Delta\theta(z;w)
    \right)
    -
    \widehat J_{\mathcal B}(\theta).
    \label{eq:app_empirical_negative_update_utility}
\end{equation}
The corresponding utility-optimal weight is
\begin{equation}
    \widehat w_{\mathcal U}^{\star}(z\mid\theta)
    \in
    \arg\max_{w\in[0,1]}
    \widehat{\mathcal U}_{\mathcal B}(z;w\mid\theta).
    \label{eq:empirical_utility_optimal_weight}
\end{equation}

We distinguish this empirical utility from its population counterpart:
\begin{equation}
    \mathcal U_\theta(z;w)
    =
    J\!\left(
        \theta+\eta w g_\theta^-(z)
    \right)
    -
    J(\theta).
    \label{eq:population_negative_update_utility}
\end{equation}
If \(\widehat J_{\mathcal B}\) is an unbiased estimator of \(J\), then
\begin{equation}
    \mathbb E_{\mathcal B}
    \left[
        \widehat{\mathcal U}_{\mathcal B}(z;w\mid\theta)
    \right]
    =
    \mathcal U_\theta(z;w).
    \label{eq:empirical_utility_unbiased}
\end{equation}
Thus, empirical utility is a noisy, evaluation-dependent estimate rather
than an intrinsic deterministic label attached to the sample.

The following result gives sufficient conditions under which the
utility-optimal effective step vanishes in the far field.

\begin{lemma}[Far-field shrinkage of the utility-optimal step]
\label{lem:far_field_utility}
Let \(D\) index learner-relative remoteness, and let \(g(D)\neq0\)
denote the corresponding negative-update direction. Write
\begin{equation}
    I(D)
    =
    \|g(D)\|_2,
    \qquad
    v(D)
    =
    \frac{g(D)}{I(D)}.
    \label{eq:far_field_direction}
\end{equation}
Assume that the local update space admits an orthogonal decomposition
\begin{equation}
    \mathcal T
    =
    \mathcal T_{\mathrm{task}}
    \oplus
    \mathcal T_{\mathrm{far}},
    \label{eq:task_far_decomposition}
\end{equation}
and let \(P_{\mathrm{task}}\) denote the orthogonal projection onto
\(\mathcal T_{\mathrm{task}}\). Suppose that:
\begin{enumerate}
    \item \(\nabla J(\theta)\in\mathcal T_{\mathrm{task}}\);

    \item the normalized far-field update becomes asymptotically
    task-orthogonal,
    \begin{equation}
        \left\|
            P_{\mathrm{task}}v(D)
        \right\|_2
        \longrightarrow 0;
        \label{eq:task_projection_vanishes}
    \end{equation}

    \item there exist \(\rho_0>0\), \(\kappa>0\), and \(D_0\) such that,
    for every \(D\geq D_0\) and every \(\rho\in[0,\rho_0]\),
    \begin{equation}
        v(D)^\top
        \nabla^2J
        \!\left(
            \theta+\rho v(D)
        \right)
        v(D)
        \leq
        -\kappa.
        \label{eq:far_field_negative_curvature}
    \end{equation}
\end{enumerate}

Define the population utility of an effective step of length \(\rho\) by
\begin{equation}
    \mathcal V_D(\rho)
    =
    J\!\left(
        \theta+\rho v(D)
    \right)
    -
    J(\theta),
    \qquad
    0\leq\rho\leq\rho_0,
    \label{eq:effective_step_utility}
\end{equation}
and let
\begin{equation}
    \rho_{\mathcal U}^{\star}(D)
    \in
    \arg\max_{\rho\in[0,\rho_0]}
    \mathcal V_D(\rho).
    \label{eq:utility_optimal_effective_step}
\end{equation}
Then
\begin{equation}
    \rho_{\mathcal U}^{\star}(D)
    \longrightarrow 0
    \qquad
    \text{as }
    D\longrightarrow\infty.
    \label{eq:optimal_step_vanishes}
\end{equation}
Moreover, if
\[
    \langle\nabla J(\theta),v(D)\rangle<0,
\]
then every nonzero local step has negative utility:
\begin{equation}
    \mathcal V_D(\rho)<0
    \qquad
    \text{for all }
    \rho\in(0,\rho_0].
    \label{eq:negative_far_field_utility}
\end{equation}
\end{lemma}

\begin{proof}
Define the first-order task alignment
\begin{equation}
    a(D)
    =
    \left\langle
        \nabla J(\theta),
        v(D)
    \right\rangle.
    \label{eq:far_field_alignment}
\end{equation}
Because
\(\nabla J(\theta)\in\mathcal T_{\mathrm{task}}\), orthogonality gives
\begin{equation}
\begin{aligned}
    |a(D)|
    &=
    \left|
        \left\langle
            \nabla J(\theta),
            P_{\mathrm{task}}v(D)
        \right\rangle
    \right| \\
    &\leq
    \left\|
        \nabla J(\theta)
    \right\|_2
    \left\|
        P_{\mathrm{task}}v(D)
    \right\|_2.
\end{aligned}
\end{equation}
Equation~\eqref{eq:task_projection_vanishes} therefore implies
\begin{equation}
    a(D)\longrightarrow0.
    \label{eq:alignment_vanishes}
\end{equation}

For any \(\rho\in[0,\rho_0]\), Taylor's theorem gives some
\(\xi\in(0,\rho)\) such that
\begin{equation}
\begin{aligned}
    \mathcal V_D(\rho)
    &=
    \rho a(D) \\
    &\quad+
    \frac{\rho^2}{2}
    v(D)^\top
    \nabla^2J
    \!\left(
        \theta+\xi v(D)
    \right)
    v(D).
\end{aligned}
\end{equation}
Using Equation~\eqref{eq:far_field_negative_curvature},
\begin{equation}
    \mathcal V_D(\rho)
    \leq
    \rho a(D)
    -
    \frac{\kappa}{2}\rho^2.
    \label{eq:utility_quadratic_upper_bound}
\end{equation}

Since \(\mathcal V_D(0)=0\), any utility-maximizing step must attain
nonnegative utility. However, whenever
\begin{equation}
    \rho
    >
    \frac{
        2[a(D)]_+
    }{
        \kappa
    },
\end{equation}
the right-hand side of
Equation~\eqref{eq:utility_quadratic_upper_bound}
is strictly negative. Hence,
\begin{equation}
    0
    \leq
    \rho_{\mathcal U}^{\star}(D)
    \leq
    \frac{
        2[a(D)]_+
    }{
        \kappa
    }.
    \label{eq:optimal_step_upper_bound}
\end{equation}
Together with Equation~\eqref{eq:alignment_vanishes}, this proves
Equation~\eqref{eq:optimal_step_vanishes}.

If \(a(D)<0\), then for every \(\rho>0\),
\begin{equation}
    \rho a(D)
    -
    \frac{\kappa}{2}\rho^2
    <
    0.
\end{equation}
Equation~\eqref{eq:utility_quadratic_upper_bound} therefore implies
\(\mathcal V_D(\rho)<0\) for every nonzero local step, proving
Equation~\eqref{eq:negative_far_field_utility}.
\end{proof}

The lemma is a sufficient mechanism rather than a universal monotonicity
claim. It applies when far-field negative updates become increasingly
misaligned with the task-improving direction and the task objective is
locally harmed along that far-field direction. The curvature condition
can represent local losses caused by parameter drift or policy-support
contraction, provided those effects reduce the evaluated objective
\(J\). If far-field growth remains aligned with \(\nabla J(\theta)\), the
conclusion need not hold.

Finally, the effective step generated by a sample weight \(w\) is
\begin{equation}
    \rho(D;w)
    =
    \eta w I(D).
    \label{eq:weight_effective_step}
\end{equation}
Consequently, the population analogue of the utility-optimal sample
weight must satisfy
\begin{equation}
    \eta
    w_{\mathcal U}^{\star}(D)
    I(D)
    =
    \rho_{\mathcal U}^{\star}(D)
    \longrightarrow0.
    \label{eq:oracle_weighted_influence_vanishes}
\end{equation}
Thus, when the raw far-field response remains nonvanishing or grows, the
utility-optimal sample weight itself must shrink so that the weighted
response vanishes.

\subsection{Proof of Proposition~\ref{prop:exp_vanishing_influence}}
\label{app:exp_vanishing_influence}

\begin{proof}
For $D>\tau$,
\begin{equation}
    \omega_{\mathrm{Exp}}(D)
    =
    \exp
    \left[
        -\frac{\lambda}{c}(D-\tau)
    \right].
\end{equation}
Hence,
\begin{equation}
\begin{aligned}
    \omega_{\mathrm{Exp}}(D) I(D)
    &\leq
    C
    \exp
    \left[
        -\frac{\lambda}{c}(D-\tau)
    \right]
    (1+D)^m \\
    &\longrightarrow
    0,
\end{aligned}
\end{equation}
because exponential decay dominates every finite-order polynomial.

The corresponding weighted parameter step is
\begin{equation}
    \Delta\theta_{\mathrm{Exp}}(D)
    =
    \eta
    \omega_{\mathrm{Exp}}(D)
    g_\theta^-(D),
\end{equation}
and therefore
\begin{equation}
    \left\|
        \Delta\theta_{\mathrm{Exp}}(D)
    \right\|_2
    =
    \eta
    \omega_{\mathrm{Exp}}(D)
    I(D)
    \longrightarrow
    0.
\end{equation}

If $\widehat J_{\mathcal B}$ is locally $L_{\mathcal B}$-Lipschitz around
$\theta$, then for all sufficiently large $D$,
\begin{equation}
\begin{aligned}
    \left|
        \widehat{\mathcal U}_{\theta,\mathcal B}
        \bigl(z;\omega_{\mathrm{Exp}}(D)\bigr)
    \right|
    &\leq
    L_{\mathcal B}
    \left\|
        \Delta\theta_{\mathrm{Exp}}(D)
    \right\|_2 \\
    &\longrightarrow
    0.
\end{aligned}
\end{equation}
This proves both claims.
\end{proof}


\subsection{Categorical One-vs-Rest Dynamics}
\label{app:categorical_self_limiting}

\begin{proof}[Proof of Theorem~\ref{thm:categorical_self_limiting}]
For a single negative atom \(i\), the direct-logit field is
\[
\dot z=-q(e_i-\pi),
\qquad
\dot z_i=-q(1-\pi_i),
\qquad
\dot z_j=q\pi_j\quad(j\ne i).
\]
Define the normalized off-target probabilities
\(\bar\pi_j=\pi_j/(1-\pi_i)\). Differentiating
\(y_i=z_i-\log\sum_{j\ne i}e^{z_j}\) gives
\begin{align}
\dot y_i
&=-q(1-\pi_i)
-q\frac{\sum_{j\ne i}\pi_j^2}{1-\pi_i} \\
&=-q(1-\pi_i)c_K(t),
\qquad
c_K(t)=1+\sum_{j\ne i}\bar\pi_j^2.
\label{eq:app_categorical_exact_ode}
\end{align}
Because \(\bar\pi\) is a distribution on \(K-1\) actions,
\[
\frac1{K-1}\le\sum_{j\ne i}\bar\pi_j^2\le1,
\]
which proves the coefficient envelope without a concentration assumption.

We next justify entry into the far half-line before using its bounds.
The log-odds is strictly decreasing. Until \(\pi_i<1/2\), monotonicity
gives \(1-\pi_i(t)\ge1-\pi_i(0)>0\), so
Equation~\eqref{eq:app_categorical_exact_ode} has a strictly negative
upper bound. The trajectory therefore crosses \(\pi_i=1/2\) in finite
time. Thereafter \(1-\pi_i\ge1/2\), so \(y_i\to-\infty\) at a linear
rate and \(\pi_i=\sigma(y_i)\) decays exponentially.

For the generic coefficient, suppose \(j^*\ne i\) is the unique largest
off-target logit initially. For every \(k\ne i,j^*\),
\[
\frac{d}{dt}(z_{j^*}-z_k)=q(\pi_{j^*}-\pi_k)>0.
\]
The initial ordering is therefore preserved. If a gap stayed bounded,
then, once \(\pi_i\to0\), maximality of \(j^*\) and the positive initial
gap would give a positive lower bound on
\(\pi_{j^*}-\pi_k\), a contradiction. Hence all such gaps diverge,
off-target mass concentrates on \(j^*\), and \(c_K(t)\to2\). Under
complete off-target symmetry, symmetry is preserved and
\(\bar\pi_j=1/(K-1)\), giving
\(c_K(t)\equiv K/(K-1)\).

Now apply the taper
\(\omega_i=\min\{1,(e^{\tau_i}\pi_i)^\beta\}\). In the far branch,
let \(U=e^{-\beta y_i}\). Since
\(\pi_i e^{-y_i}=1-\pi_i\),
\begin{align}
\dot U
&=-\beta e^{-\beta y_i}\dot y_i \\
&=\beta q e^{\beta\tau_i}
c_K(t)(1-\pi_i)^{\beta+1}.
\label{eq:app_categorical_tapered_u}
\end{align}
After the finite crossing above, the right-hand side is bounded above
and below by positive constants. Thus \(U=\Theta(t)\), which is
equivalent to
\[
y_i(t)=-\beta^{-1}\log t+O(1),
\qquad
\pi_i(t)=\Theta(t^{-1/\beta}).
\]
\end{proof}

\subsection{Categorical Restoring-Force Balance}
\label{app:categorical_restoring_balance}

Write the one-dimensional restored dynamics as
\[
\dot y=F(y)=E(y)-S_{\mathrm{tap}}(y),
\]
where \(S_{\mathrm{tap}}>0\) is downward suppression and \(E>0\) is
upward restoration.

\begin{proposition}[Finite restoring crossing]
\label{prop:categorical_restoring_crossing}
Assume \(E\) and \(S_{\mathrm{tap}}\) are continuous and their zeros in
\([y_L,y_R]\) are isolated. If
\(F(y_L)>0\) and \(F(y_R)<0\), then the interval contains a zero at
which the phase line has a positive-to-negative crossing. Such a
crossing is locally asymptotically stable; if it is hyperbolic, its
condition is
\[
F'(y^*)<0
\quad\Longleftrightarrow\quad
[S_{\mathrm{tap}}-E]'(y^*)>0.
\]
\end{proposition}

\begin{proof}
The intermediate value theorem gives at least one zero. Because the
zeros are isolated on a compact interval, only finitely many occur; as
the endpoint signs differ, at least one has positive sign immediately
to its left and negative sign immediately to its right. The flow then
points toward that zero from both sides, proving phase-line stability. At a
hyperbolic crossing this sign change is equivalent to \(F'(y^*)<0\),
and the displayed derivative identity follows from \(F=E-S_{\mathrm{tap}}\).
\end{proof}

In the binary or symmetric one-vs-rest tail, the tapered suppression
has order \(e^{\beta y}\), whereas entropy restoration has order
\(|y|e^y\) as \(y\to-\infty\). Hence \(\beta\ge1\) supplies the
required far-tail ordering (with the \(|y|\) factor breaking the tie at
\(\beta=1\)), but it does not by itself verify the finite phase-line
conditions of Proposition~\ref{prop:categorical_restoring_crossing}.

\section{Distributional Reweighting Interpretations}
\label{app:distributional_reweighting}

\subsection{Finite-Measure I-Divergence Attenuation}
\label{app:finite_measure_variational}

\begin{proof}[Proof of Proposition~\ref{prop:finite_measure_variational}]
Use the convention
\(\phi(0)=1\) for
\(\phi(w)=w\log w-w+1\). For an active constraint, the Lagrangian with
multiplier \(\kappa>0\) is
\[
\mathcal L(w,\kappa)
=\int\{wh+\kappa\phi(w)\}\,d\xi^- -\kappa\delta.
\]
Pointwise stationarity on \(0<w\le1\) gives
\(h+\kappa\log w=0\), hence
\[
w_\kappa(z)=e^{-h(z)/\kappa}.
\]
For \(h=0\), this equals one; for \(h>0\), it lies strictly between
zero and one. Define
\(\mathcal D(\kappa)=\int\phi(w_\kappa)d\xi^-\). Dominated convergence gives
\[
\lim_{\kappa\to\infty}\mathcal D(\kappa)=0,
\qquad
\lim_{\kappa\downarrow0}\mathcal D(\kappa)
=\xi^-(\{h>0\})=\delta_0.
\]
On a set of positive \(h\), increasing \(\kappa\) strictly increases
\(w_\kappa\) toward one and strictly decreases \(\phi(w_\kappa)\).
Thus \(\mathcal D(\kappa)\) is continuous and strictly decreasing from
\(\delta_0\) to zero, proving existence and uniqueness for
\(0<\delta<\delta_0\).

At \(\delta=\delta_0\), the limit is
\(w_0=\mathbf1\{h=0\}\), which has zero objective and divergence
\(\delta_0\). For \(\delta>\delta_0\), the same zero-objective point is
strictly feasible, so the constraint need not be active and there is no
finite one-to-one multiplier relation. Finally,
\(\int\phi(0)d\xi^-=|\xi^-|\), so \(w\equiv0\) becomes feasible only at the later
endpoint \(\delta\ge|\xi^-|\).
\end{proof}

This problem derives learner-relative exponential attenuation. It is
distinct from quality-based hard selection: bo